% 由 scripts/publication/tables.py 自动生成，请勿手改。
% 需要宏包：booktabs（\usepackage{booktabs}）
\begin{table}[htbp]
  \centering
  \caption{Step 13 补修后主结果：固定端点六方案（256/512）}
  \label{tab:step13-main}
  \begin{tabular}{lrrrrrr}
    \toprule
    方案 & 平均损耗 (dB) & P95 (dB) & 最大损耗 (dB) & 直线 (dB) & 弯曲 (dB) & 交叉 (dB) \\
    \midrule
    A & 5.2786 & 6.1318 & 6.3333 & 0.6249 & 4.5228 & 0.1308 \\
    R5U & 5.2778 & 6.1310 & 6.3222 & 0.6271 & 4.5228 & 0.1278 \\
    D0 & 4.3815 & 5.1735 & 5.3638 & 0.6088 & 3.6375 & 0.1352 \\
    D56 & 4.3393 & 5.1587 & 5.3462 & 0.6181 & 3.5889 & 0.1323 \\
    F5 & 3.4981 & 5.5981 & 5.9980 & 0.5857 & 2.6525 & 0.2598 \\
    F56 & 2.9703 & 4.6545 & 5.0783 & 0.5842 & 2.1202 & 0.2659 \\
    \midrule
    A & 5.5162 & 6.3510 & 6.5773 & 0.6355 & 4.5587 & 0.3219 \\
    R5U & 5.5167 & 6.3484 & 6.5780 & 0.6353 & 4.5587 & 0.3227 \\
    D0 & 4.5747 & 5.3886 & 5.6207 & 0.6204 & 3.6190 & 0.3353 \\
    D56 & 4.5607 & 5.3886 & 5.6217 & 0.6225 & 3.6040 & 0.3343 \\
    F5 & 3.8261 & 5.8453 & 6.4459 & 0.5969 & 2.6664 & 0.5628 \\
    F56 & 3.2947 & 4.9004 & 5.4995 & 0.5913 & 2.1384 & 0.5650 \\
    \bottomrule
  \end{tabular}
  \par\vspace{2pt}
  \begin{flushleft}\footnotesize
  A=原版几何（端点来源）；R5U=固定端点+U 型+R5；D0=优先 R6 回退 R5；D56=自适应 R5/R6；F5/F56=自由弯角（半径 R5 / R5+R6）。 \\
  全部方案端点冻结（endpoint\_changes=0），D56/F56 为本次补修后（256 的 F56 已重布历史重合对 \#7/\#31）。 \\
  自由弯角方案以更高的交叉损耗换取更低的弯曲损耗（见分解列），最优值加粗仅在同列六方案内。 \\
  损耗为解析物理统计口径（交叉事件按每根波导各计一次）；与原版 calc\_index 统计口径不同，两者不混用。 \\
  \end{flushleft}
\end{table}
